\documentclass[10pt]{beamer}

\usetheme[progressbar=frametitle]{metropolis}
\usecolortheme{aggie}

\usepackage{appendixnumberbeamer}

\usepackage{booktabs}
\usepackage[scale=2]{ccicons}



\usepackage{pgfplots}
\usepgfplotslibrary{dateplot}

\usepackage{xspace}
\newcommand{\themename}{\textbf{\textsc{metropolis}}\xspace}

\newcommand{\incfig}[1]{%
\center
\def\svgwidth{0.9\columnwidth}
\import{./figures/}{#1.pdf_tex}
}

\title{Introduction to sheaves and schemes}
\subtitle{Final presentation for the subject of algebraic geomety}
% \date{\today}
\date{}
\author{Abel Doñate Muñoz}
\institute{Universitat Politècnica de Catalunya}
% \titlegraphic{\hfill\includegraphics[height=1.5cm]{logo.pdf}}

\begin{document}

\maketitle

\begin{frame}{Table of contents}
  \setbeamertemplate{section in toc}[sections numbered]
  \tableofcontents%[hideallsubsections]
\end{frame}

\section[Intro]{Introduction}

\begin{frame}[fragile]{Definición de prehaz}
Sea $X$ un espacio topológico. Un haz es una apicación . 

%\incfig{sheaf}

\end{frame}


\begin{frame}[fragile]{Definición de haz}
Sea $X$ un espacio topológico. Un haz es una apicación . 
\end{frame}


\begin{frame}[fragile]{Morfismos de haces}
A
\end{frame}


\begin{frame}[fragile]{Núcleo e imagen de morfismos de haces}
A
\end{frame}


\begin{frame}[fragile]{Stalk}
A
\end{frame}


\begin{frame}[fragile]{Ejemplos de haces y stalks}
A
\end{frame}








\begin{frame}[allowframebreaks]{References}


\end{frame}

\end{document}
